%% 1-0 序章　かずを つかうと、なにが いいの？（4ページ、p.10〜13）
\makeatletter
\gdef\bk@partlabel{\R{第}{だい}1\R{巻}{かん}}\gdef\bk@partname{かぞえる・はかる}\gdef\bk@tabpos{0}
\makeatother
\bkchapter{1}{0}{かずを つかうと、\\なにが いいの？}{\bklevel{1}}{かずを つかうと、なにが いいの？}

\noindent あしたは えんそく。 ソラと ミナが、 おやつの あめを みせあって います。

\begin{bkfigure}
\begin{tikzpicture}[x=1mm, y=1mm]
  \fill[bkYellowBg, rounded corners=4mm] (0,0) rectangle (132,62);
  % ソラ
  \bkChar{sora}{(19,4)}{34mm}
  \node[rectangle callout, callout absolute pointer={(30,34)}, draw=bkTheme, line width=0.8pt, fill=white,
    rounded corners=2.5mm, inner xsep=2.5mm, inner ysep=1.6mm, font=\small, align=left, anchor=south west] at (4,42) {あめ、 たくさん\\もって きたよ！};
  \fill[white, rounded corners=2mm] (36,6) rectangle (62,30);
  \foreach \i in {0,...,7} {\pgfmathtruncatemacro\cx{mod(\i,3)}\pgfmathtruncatemacro\cy{div(\i,3)}
    \bkcandy{(41+\cx*8,25-\cy*7)}{1.6}{bkRoute4}}
  % ミナ
  \bkChar{mina}{(113,4)}{34mm}
  \node[rectangle callout, callout absolute pointer={(102,34)}, draw=bkTheme, line width=0.8pt, fill=white,
    rounded corners=2.5mm, inner xsep=2.5mm, inner ysep=1.6mm, font=\small, align=left, anchor=south east] at (128,42) {わたしも\\たくさん！};
  \fill[white, rounded corners=2mm] (66,6) rectangle (96,30);
  \foreach \i in {0,...,11} {\pgfmathtruncatemacro\cx{mod(\i,4)}\pgfmathtruncatemacro\cy{div(\i,4)}
    \bkcandy{(70.5+\cx*7,25-\cy*7)}{1.6}{bkRoute6}}
  \node[font=\sffamily\bfseries\Large, text=bkTheme] at (66,47) {？};
\end{tikzpicture}
\end{bkfigure}

\noindent ふたりとも 「たくさん」。 でも、 どっちが おおいのかな？\\
「たくさん」と いう ことばだけでは、 くらべられません。

%% ---------------- かずに すると、くらべられる（第11回：ペアと かぞえるの 2とおりで みせる） ----------------
\bksubsection{1}{かずに すると、 くらべられる}

\noindent あめを 1こずつ ならべて、 ソラの あめと ミナの あめを 1つずつ ペアに して みましょう。

\begin{bkfigure}
\begin{tikzpicture}[x=1mm, y=1mm]
  \fill[bkYellowBg, rounded corners=3mm] (-2,-8) rectangle (130,24);
  \node[anchor=west, font=\sffamily\bfseries\small, text=bkInk] at (0,15) {ソラ};
  \node[anchor=west, font=\sffamily\bfseries\small, text=bkInk] at (0,0) {ミナ};
  \foreach \i in {0,...,7} {\bkcandy{(24+\i*9,15)}{1.6}{bkRoute4}}
  \foreach \i in {0,...,11} {\bkcandy{(24+\i*9,0)}{1.6}{bkRoute6}}
  \foreach \i in {0,...,7} {\draw[bkGray, line width=0.8pt] (24+\i*9,12.2) -- (24+\i*9,2.8);}
  \draw[bkTheme, line width=1.2pt] (110.5,0) ellipse (18 and 4.8);
  \node[font=\sffamily\bfseries\small, text=bkTheme] at (110.5,15) {あまった 4こ};
  \node[font=\sffamily\footnotesize, text=bkGray!80!black] at (55.5,-5.6) {せんで むすんだ ペア};
\end{tikzpicture}
\end{bkfigure}

\noindent ミナの あめが 4こ あまりました。 あまった ほうが おおいね。

\noindent こんどは、 1こずつ ゆびで さしながら かぞえて みましょう。 ソラは \textbf{8こ}、 ミナは \textbf{12こ}。\\
かぞえる ときは、 1、 2、 3、 …と、 じゅんばんに いいますね。 8の あとに 9、 10、 11、 12と つづきます。 あとから でて くる かずの ほうが、 おおきい かずです。 だから、 12は 8より おおきい。 ならべなくても、 かずを きくだけで 「ミナの ほうが おおい」と わかります。

\begin{taisetsu}
かずに すると、 \textbf{くらべられる}。
\end{taisetsu}

\noindent かずの かぞえかたと くらべかたは、 1-1で ゆっくり やります。 12のような 10より おおきい かずは、 1-2で まなびます。

%% ---------------- p.11 つたえられる・さきが わかる ----------------
\bksubsection{2}{かずに すると、 \R{人}{ひと}に つたえられる}

\noindent\begin{minipage}[c]{0.62\linewidth}
「リンゴを たくさん かって きてね」と たのまれたら、 なんこ かえば いいのか こまって しまいます。 さっきの 「たくさん」も、 ソラは 8こ、 ミナは 12こでした。 「たくさん」は、 \R{人}{ひと}に よって ちがうのです。

「リンゴを 6こ かって きてね」なら、 だれが いっても、 おなじ かずだけ かって こられます。
\end{minipage}\hfill
\begin{minipage}[c]{0.35\linewidth}\centering
\begin{tikzpicture}[x=1mm, y=1mm]
  \fill[white, draw=bkLine, line width=0.8pt, rounded corners=3mm] (0,0) rectangle (46,34);
  \bkplate{(23,9)}{38}
  \foreach \x/\y in {10/12,18/14,26/14.5,34/13.5,14/20.5,23/21.5} {\bkapple{(\x,\y)}{3.6}{bkAppleRed}}
  \node[font=\sffamily\bfseries\small, text=bkTheme] at (37,27) {6こ};
\end{tikzpicture}
\end{minipage}

\bksubsection{3}{かずに すると、 この さきが わかる}

\noindent\begin{minipage}[c]{0.6\linewidth}
\R{水}{みず}そうに バケツで \R{水}{みず}を いれます。 バケツ 1ぱいで、 \R{水}{みず}は \R{線}{せん}（\R{目}{め}もり） 1つ\R{分}{ぶん} あがります。

いまは \R{目}{め}もり 6。 いっぱいは \R{目}{め}もり 10です。 あと なんばいで いっぱいに なるかな？

いまの 6の つぎから、 7、 8、 9、 10と \R{目}{め}もりを かぞえると 4つ。 バケツ 1ぱいで \R{目}{め}もり 1つ だから、 あと 4はいです。 \R{水}{みず}を いれる まえに、 かずで わかりました。
\end{minipage}\hfill
\begin{minipage}[c]{0.37\linewidth}\centering
\begin{tikzpicture}[x=1mm, y=1mm]
  % 水そう
  \fill[bkBlue!25] (2,0.6) rectangle (34,30);
  \draw[bkInk!70, line width=1.2pt, line join=round] (1.4,50) -- (1.4,0) -- (34.6,0) -- (34.6,50);
  \foreach \k in {1,...,10} {\draw[bkInk, line width=0.6pt] (34.6,\k*5) -- (30.6,\k*5);
    \node[anchor=west, font=\sffamily\bfseries\footnotesize, text=bkInk] at (35.4,\k*5) {\k};}
  \draw[bkRed, line width=1pt, dash pattern=on 2pt off 1.5pt] (1.4,50) -- (34.6,50);
  \node[anchor=south, font=\sffamily\bfseries\footnotesize, text=bkRed] at (18,50.5) {いっぱい};
  \foreach \k in {7,...,10} {\draw[bkOrange, line width=1.1pt, -{Stealth[length=3pt]}] (25,{\k*5-4.3}) -- (25,{\k*5-0.7});}
  \node[anchor=east, font=\sffamily\bfseries\footnotesize, text=bkOrange!80!black, align=right] at (23,42.5) {あと\\4はい};
  % バケツ
  \begin{scope}[shift={(-2,60)}, rotate=-35]
    \fill[bkGray!50] (0,0) -- (12,0) -- (10.5,-10) -- (1.5,-10) -- cycle;
    \draw[bkInk!70, line width=0.8pt] (0,0) .. controls (2,5) and (10,5) .. (12,0);
  \end{scope}
\end{tikzpicture}
\end{minipage}

\par\vspace{8pt}
\noindent カレンダーでも、 おなじように さきが わかります。 きょうは \R{5日}{いつか}。 えんそくは \R{8日}{ようか}です。

\begin{bkfigure}
\begin{tikzpicture}[x=1mm, y=1mm]
  \fill[bkTheme!6, rounded corners=3mm] (0,-9) rectangle (132,31);
  \foreach \d [count=\i from 0] in {3,4,5,6,7,8,9} {
    \filldraw[fill=white, draw=bkGray, line width=0.7pt] (8+\i*17,0) rectangle ++(15,15);
    \node[font=\sffamily\bfseries\large, text=bkInk] at (15.5+\i*17,7.5) {\d};}
  \draw[bkBlue, line width=1.3pt] (49.5,7.5) circle (6);
  \node[font=\sffamily\bfseries\footnotesize, text=bkBlue!80!black] at (49.5,-4.5) {きょう};
  \node[font=\sffamily\bfseries\footnotesize, text=bkTheme] at (100.5,-4.5) {えんそく};
  \draw[bkTheme, line width=1.3pt] (100.5,7.5) circle (6);
  \foreach \i in {0,1,2} {\draw[bkOrange, line width=1.3pt, -{Stealth[length=5pt]}] (51+\i*17,16.5) .. controls (55+\i*17,26) and (61+\i*17,26) .. (65+\i*17,16.5);
    \node[circle, fill=bkOrange, text=white, font=\sffamily\bfseries\footnotesize, minimum size=4.6mm, inner sep=0pt] at (58+\i*17,27) {\the\numexpr\i+1\relax};}
\end{tikzpicture}
\end{bkfigure}

\noindent 5の つぎから、 6、 7、 8と かぞえると 3つ。 あと 3かい ねると、 えんそくの \R{日}{ひ}です。 かずと じゅんばんが わかって いれば、 その \R{日}{ひ}が くる まえに わかります。

\clearpage
%% ---------------- きかいにも つたえられる（第11回で追加） ----------------
\bksubsection{4}{かずなら、 きかいにも つたえられる}

\noindent きかいは、 「たくさん」や 「すこし」が わかりません。 でも、 かずで つたえると、 そのとおりに うごきます。

\begin{bkfigure}
\begin{tikzpicture}[x=1mm, y=1mm,
  btn/.style={circle, draw=bkGray, line width=0.8pt, fill=white, minimum size=7.5mm, inner sep=0pt, font=\sffamily\bfseries\small}]
  % エレベーターの ボタン
  \fill[white, draw=bkLine, line width=0.8pt, rounded corners=3mm] (0,0) rectangle (62,58);
  \fill[bkGray!25, rounded corners=2mm] (16,10) rectangle (46,54);
  \fill[bkInk!85, rounded corners=1mm] (21,44) rectangle (41,51);
  \node[font=\sffamily\bfseries\small, text=bkYellow] at (31,47.5) {5かい};
  \foreach \n/\x/\y in {1/24/15,2/38/15,3/24/25,4/38/25,6/38/35} {\node[btn] at (\x,\y) {\n};}
  \node[btn, fill=bkYellow!60, draw=bkOrange] at (24,35) {5};
  \node[font=\sffamily\bfseries\small] at (31,4.5) {エレベーター};
  % でんしレンジ
  \begin{scope}[shift={(70,0)}]
  \fill[white, draw=bkLine, line width=0.8pt, rounded corners=3mm] (0,0) rectangle (62,58);
  \fill[bkGray!35, rounded corners=2mm] (5,14) rectangle (57,46);
  \fill[bkInk!80, rounded corners=1mm] (9,18) rectangle (40,42);
  \fill[bkGray!15, rounded corners=1mm] (43,18) rectangle (54,42);
  \fill[bkGreen!25] (44.5,34) rectangle (52.5,40);
  \node[font=\sffamily\bfseries\footnotesize, text=bkInk] at (48.5,37) {30};
  \foreach \y in {23,29} {\fill[bkGray!70] (48.5,\y) circle (1.8);}
  \node[font=\sffamily\bfseries\footnotesize, text=bkInk] at (31,52) {30びょう あたためる};
  \node[font=\sffamily\bfseries\small] at (31,4.5) {でんしレンジ};
  \end{scope}
\end{tikzpicture}
\end{bkfigure}

\begin{itemize}[label={\color{bkPurple}●}, leftmargin=1.4\zw, itemsep=0pt, topsep=0pt]
\item エレベーターに 「うえの ほう」と いっても、 どこで とまれば いいか わかりません。 「5」の ボタンを おすと、 5かいで とまります。
\item でんしレンジは、 「30」と かずで いれると、 30びょうだけ あたためます。
\end{itemize}

\noindent ロボットも ロケットも おなじです。 \R{人}{ひと}が かずで つたえるから、 きまった とおりに うごけるのです。

\begin{taisetsu}
かずに すると、 \textbf{くらべられる}。 \textbf{つたえられる}。\\
\textbf{この さきが わかる}。 \textbf{きかいも うごかせる}。
\end{taisetsu}


\clearpage
%% ---------------- p.12 この本で できるように なる こと ----------------
\bksection{goal}{この \R{本}{ほん}で できるように なる こと}

\noindent この \R{本}{ほん}は、 4つの 「ぶ」に わかれて います。

\begin{bkfigure}
\begin{tikzpicture}[x=1mm, y=1mm,
  pb/.style={draw=bkVol, line width=1.2pt, fill=white, rounded corners=3mm, minimum width=63mm, minimum height=50mm}]
  \foreach \i/\x/\y in {1/0/53,2/69/53,3/0/0,4/69/0} {\node[pb, anchor=south west] (p\i) at (\x,\y) {};}
  \foreach \i/\x/\y/\t/\c/\pg in {
      1/0/53/{かずの はじまり}/{かぞえる・ くらべる・ たす・ ひく}/{1-1〜1-7},
      2/69/53/はかる/{ながさ・ かさ・ おもさ・ \R{時間}{じ|かん}を はかる}/{1-8〜1-11},
      3/0/0/{かけ\R{算}{ざん}と わり\R{算}{ざん}}/{\R{同}{おな}じ かずの まとまりで かぞえる・ \R{分}{わ}ける}/{1-12〜1-16},
      4/69/0/{\R{表}{ひょう}・グラフ・かたち}/{\R{表}{ひょう}や グラフに する・ かたちと \R{広}{ひろ}さ}/{1-17〜1-21}} {
    \fill[bkVol, rounded corners=2mm] (\x+3,\y+40) rectangle (\x+17,\y+47);
    \node[text=bkOn, font=\sffamily\bfseries\footnotesize] at (\x+10,\y+43.2) {\R{第}{だい}\i\R{部}{ぶ}};
    \node[anchor=west, font=\sffamily\bfseries\small, text=bkDeep] at (\x+19,\y+43.2) {\t};
    \node[anchor=north west, font=\footnotesize, text=bkInk, text width=57mm, align=left, inner sep=0pt] at (\x+4,\y+36.5) {\c};
    \node[anchor=south east, font=\sffamily\bfseries\footnotesize, text=bkGray] at (\x+61,\y+1.5) {\pg};}% 第11回：ページ数が かわるので、ページ番号ではなく 章の はんいを かく
  % 第1部：リンゴと 10のわく
  \foreach \x in {7,14,21} {\bkapple{(\x,65)}{3}{bkAppleRed}}
  \node[font=\sffamily\bfseries\small, text=bkTheme] at (28,65) {＋};
  \foreach \x in {35,42} {\bkapple{(\x,65)}{3}{bkAppleGreen}}
  \bktenframe{(5,54.5)}{3.2}{3}{bkAppleRed}{2}{bkAppleGreen}
  % 第2部：ものさし・時計・ます
  \begin{scope}[shift={(69,53)}]
    \fill[bkYellowBg, draw=bkYellow!80!black, line width=0.6pt, rounded corners=0.8mm] (4,14) rectangle (30,20);
    \foreach \k in {0,...,10} \draw[bkInk, line width=0.4pt] (5.5+\k*2.3,20) -- (5.5+\k*2.3,{\k==5 ? 17 : 18.4});
    \bkclock*{(38,15)}{7}{3}{0}
    \bkmasu{(10,1.5)}{7}{0.6}{0}
  \end{scope}
  % 第3部：おはじきの ならび
  \foreach \i in {0,...,3} \foreach \j in {0,1,2} {\bkohajiki{(8+\i*7,{6+\j*6})}{2.3}{bkRoute4}}
  \node[font=\sffamily\bfseries\small, text=bkInk] at (42,12) {3 × 4};
  % 第4部：ぼうグラフと まる
  \begin{scope}[shift={(69,0)}]
    \draw[bkInk!70, line width=0.6pt] (5,4) -- (34,4) (5,4) -- (5,22);
    \foreach \h/\x/\c in {10/9/bkRoute2,16/17/bkRoute6,7/25/bkRoute3} {\fill[\c] (\x,4) rectangle ++(5,\h);}
    \draw[bkRoute4, line width=1.2pt] (44,15) circle (7);
    \draw[bkRoute4, line width=0.8pt] (44,15) -- ++(30:7);
    \fill[bkRoute4] (44,15) circle (0.8);
  \end{scope}
\end{tikzpicture}
\end{bkfigure}

\noindent えの なかの 「＋」や 「×」は、 これから この \R{本}{ほん}で まなぶ しるしだよ。 いまは わからなくても だいじょうぶ。

\noindent どの \R{章}{しょう}も 「はてな？」から はじまります。 さいごには、 その ふしぎに かずで こたえられるように なるよ。

\noindent \R{章}{しょう}の おわりの 「ふりかえり」で、 できるように なった ことに チェックを つけよう。 \R{本}{ほん}の うしろの 「チェックひょう」にも、 まとめて あります。

\clearpage
%% ---------------- p.13 3人の かいわ ----------------
\bksection{reidai}{3\R{人}{にん}の なかまと いっしょに}

\bkfukidashi{sora}{かずって、 かぞえる ためだけの もの だと おもって いたよ。}
\bkfukidashi{mina}{はかる ときも、 くらべる ときも、 かずを つかうんだね。}
\bkfukidashi{delta}{そうだよ。 ロケットも ロボットも、 かずで うごいて いるんだ。 1から いっしょに はじめよう！}

\begin{bkfigure}
\begin{tikzpicture}[x=1mm, y=1mm]
  \fill[bkTheme!6, rounded corners=4mm] (0,-1) rectangle (132,40);
  % 1巻から9巻への道
  \draw[bkTheme, line width=2.4pt, line cap=round] (10,12) .. controls (40,12) and (50,30) .. (80,28) .. controls (100,27) and (110,20) .. (122,26);
  \foreach \x/\y/\k in {10/12/1,27/13.5/2,42/19/3,57/25.5/4,72/28.3/5,87/27.6/6,101/25.3/7,113/22.6/8,122/26/9} {
    \fill[white, draw=bkVol\k, line width=1.4pt] (\x,\y) circle (2.8);
    \node[font=\sffamily\bfseries\scriptsize, text=bkVol\k!85!black] at (\x,\y) {\k};}
  \node[fill=bkTheme, text=white, rounded corners=1mm, font=\sffamily\bfseries\footnotesize, inner sep=1.5pt] at (10,20) {いま ここ};
  \node[anchor=north west, font=\footnotesize, text=bkInk, align=left] at (4,8) {1\R{巻}{かん}：かぞえる・ はかる};
  \node[anchor=north east, font=\footnotesize, text=bkInk, align=right] at (130,18) {9\R{巻}{かん}：ロボットや\\ロケットを うごかす すうがく};
\end{tikzpicture}
\end{bkfigure}

\noindent この \R{本}{ほん}は、 9さつ つづく \R{本}{ほん}の 1さつめです。 1さつずつ すすむと、 さいごには、 ロボットや ロケットを うごかす すうがくまで とどきます。

\noindent この \R{本}{ほん}の つかいかたは、 8ページを みてね。 わからない ところが あっても だいじょうぶ。 ゆっくり、 いっしょに すすもう。
